%Esta plantilla para Tesis de la Universidad de Santiago de Chile (USACH) está basada en la tesis de pregrado "Infinito: propuesta de formalización" de Sergio Rubio Pizzorno (www.zergiorubio.org) la cual se puede encontrar en el Repositorio Digital de la USACH (http://repositorio.usach.cl/ se puede buscar por el nombre de la tesis -copiar y pegar-)
%Tal documento se puede mirar a modo de ejemplo para responder preguntas típicas de formato de la tesis, como el párrafo de derechos de autor y otros.
%La página http://www.zergiorubio.org/?p=3908 está dedicada a los detalles del presente Template.
%En la página http://www.zergiorubio.org/?p=3279 hay una publicación sobre la plantilla, donde se explica el uso de paquetes claves y los argumentos que respaldan la elección para utilizarlos en la plantilla.


\documentclass[12pt]{book}
\usepackage[utf8]{inputenc}
\usepackage[T1]{fontenc}
\usepackage[spanish]{babel}
\usepackage[authoryear,datebegin]{flexbib}		%Uso de flexbib para la bibliografía.
\usepackage{graphics}
\usepackage{pst-all,pstricks-add,pst-math}
\usepackage{tikz}
\usepackage{float}
\usepackage{textcomp}
\usepackage{dsfont}
\usepackage{booktabs}		%Mejora el uso de las tablas (elimina las líneas verticales)
\usepackage{makeidx}		%Incluir Índice Analítico
\usepackage[colorlinks=true, pdfborder={0 100 100}, linkcolor=orange, linktoc=page, citecolor=Green]{hyperref}
\usepackage{microtype}		%Arregla el aspecto gráfico del documento (incluye una mejor separación entre letras y palabras)
\usepackage{dsfont,amssymb,amsfonts,latexsym,cancel}

\definecolor{orange}{RGB}{250,120,0}
\definecolor{Green}{RGB}{50,200,100}

\newcommand{\N}{\mathds{N}}	%Redefinición símbolo números Naturales
\newcommand{\Z}{\mathds{Z}}	%Redefinición símbolo números Enteros
\newcommand{\Q}{\mathds{Q}}	%Redefinición símbolo números Racionales
\newcommand{\I}{\mathds{I}}	%Redefinición símbolo números Irracionales
\newcommand{\R}{\mathds{R}}	%Redefinición símbolo números Reales
\newcommand{\p}{\ldotp}		%Redefinición símbolo 
\newtheorem{teorema}{Teorema}
\newtheorem{lema}{Lema}
\newtheorem{definicion}{Definición}
\newtheorem{ax}{Axioma}
\newtheorem{corolario}{Corolario}
\renewcommand{\rmdefault}{phv} % Arial
\renewcommand{\sfdefault}{phv} % Arial
\makeindex

\textheight=21.5cm    %alto del texto
\textwidth=15cm     %ancho del texto	
\topmargin=0cm     %margen superior
\oddsidemargin=1.5cm  %margen izquierdo (por defecto queda a 1.5 cm)

\usepackage{multirow}
\setcounter{tocdepth}{4} 	%Inserta subsubsecciones y párrafos en el indice
\usepackage{cleveref}		%Inserta la palabra correspondiente al referenciar un objeto (Figura, Tabla, etc)
\crefname{equation}{ecuación}{ecuaciones}
\crefname{table}{tabla}{tablas}
\crefname{figure}{figura}{figuras}
\crefname{section}{sección}{Sección}
\crefname{chapter}{capítulo}{Capítulo}

\graphicspath{{Imagenes/}} 	%La carpeta de dibujos es Imagenes
 
\includeonly{Capitulo_02}	%Incluye sólo los archivos que se especifican como parámetro. Se mantiene numeración global del documento.



\begin{document}
\renewcommand{\contentsname}{Tabla de contenidos}	%Cambiar nombre de Índice por Tabla de contenidos
\renewcommand{\listfigurename}{Índice de Ilustraciones}	%Cambiar nombre de Índice de Figuras por Índice de Ilustraciones
\renewcommand{\indexname}{Índice analítico}		%Cambiar nombre de Índice alfabético -> Índice analítico
% \renewcommand{\figurename}{Índice de Ilustraciones}	%Cambiar nombre de Índice de Figuras por Índice de Ilustraciones
\renewcommand{\listtablename}{Índice de tablas}
\renewcommand{\tablename}{Tabla}

%%%%%%%%%%%%%%%%%%%%%%%%%%%%%%%%%%
    %PÁGINAS PRELIMINARES
%%%%%%%%%%%%%%%%%%%%%%%%%%%%%%%%%%
%(1) PORTADA 			[L]
%(2) DERECHOS DE AUTOR 		[L]
%(3) HOJA DE CALIFICACIONES	[]
%(4) DEDICATORIA(OPTATIVA)	[]
%(5) AGRADECIMIENTOS(OPTATIVA)	[]
%(6) TABLA DE CONTENIDOS	[]
%(7) ÍNDICE DE TABLAS		[]
%(8) ÍNDICE DE ILUSTRACIONES	[]
%(9) RESUMEN			[]
%%%%%%%%%%%%%%%%%%%%%%%%%%%%%%%%%%

%%%%%%%%%%%%%%%%%%%%%%%%%%%%%%%%%%
%(1) PORTADA			[L]
%%%%%%%%%%%%%%%%%%%%%%%%%%%%%%%%%%
\begin{titlepage}
  \begin{flushright}
    \begin{tabular}{cl}
      &\multirow{4}{*}{\hspace{0cm}\includegraphics[width=2cm]{logoBN.png}}\\
      \large{\textbf{UNIVERSIDAD DE SANTIAGO DE CHILE}}&  \\
      \textbf{FACULTAD ...}&\\
      \textbf{DEPARTAMENTO DE ...}& \\
    \end{tabular}
  \end{flushright}

\vspace*{3cm}

  \begin{center}
    \large{\textbf{NOMBRE TESIS}}\\ %\vspace*{2cm}
     %\textbf{Estudio del concepto y su presencia en la Enseñanza de la Matemática} \vspace*{2cm}
  \end{center}
  \begin{center}
    \textbf{AUTOR}
  \end{center}

\vspace*{1cm}

  \begin{flushright}
    \begin{tabular}{ll}
      \textbf{Profesor Guía:}& NOMBRE PROF. 1\\
      & NOMBRE PROF. 2\\
    \end{tabular}
\\\vspace{3cm}

  \begin{center}
    \textbf{Trabajo de Titulación presentado en conformidad a los
    requisitos para obtener el grado de ...}
  \end{center}
  \end{flushright}

\vspace{2cm}

  \begin{center}
    \textbf{Santiago -  Chile\\AÑO}
  \end{center}
\end{titlepage}

\newpage
\mbox{}
\thispagestyle{empty} % para que no se numere esta página
\pagenumbering{roman}% para comenzar la numeración de paginas en números romanos

%%%%%%%%%%%%%%%%%%%%%%%%%%%%%%%%%%
%(2) DERECHOS DE AUTOR		[L]
%%%%%%%%%%%%%%%%%%%%%%%%%%%%%%%%%%
\vspace*{17cm}
\begin{center}
\textbf{ año, Autor\\
Derechos de autor}
\end{center}

%%%%%%%%%%%%%%%%%%%%%%%%%%%%%%%%%%
%(3) HOJA DE CALIFICACIONES	[]
%%%%%%%%%%%%%%%%%%%%%%%%%%%%%%%%%%

%%%%%%%%%%%%%%%%%%%%%%%%%%%%%%%%%%
%(4) DEDICATORIA(OPTATIVA)	[]
%%%%%%%%%%%%%%%%%%%%%%%%%%%%%%%%%%
\newpage
\chapter*{}
\begin{flushright}
\textit{Dedicatoria}
\end{flushright}
\newpage

%%%%%%%%%%%%%%%%%%%%%%%%%%%%%%%%%%
%(5) AGRADECIMIENTOS(OPTATIVA)	[]
%%%%%%%%%%%%%%%%%%%%%%%%%%%%%%%%%%
\chapter*{Agradecimientos} % *si no queremos que añada la palabra "Capitulo"
% \addcontentsline{toc}{chapter}{Agradecimientos} % si queremos que aparezca en el índice
%TEXO


%%%%%%%%%%%%%%%%%%%%%%%%%%%%%%%%%%
%(6) TABLA DE CONTENIDOS 	[]
%%%%%%%%%%%%%%%%%%%%%%%%%%%%%%%%%%
\newpage
\tableofcontents
\newpage
\newrgbcolor{xdxdff}{0.49 0.49 1}
\newrgbcolor{ffqqtt}{1 0 0.2}
\psset{xunit=1.0cm,yunit=1.0cm,algebraic=true,dotstyle=*,dotsize=3pt
0,linewidth=0.8pt,arrowsize=3pt 2,arrowinset=0.25}

%%%%%%%%%%%%%%%%%%%%%%%%%%%%%%%%%%
%(7) ÍNDICE DE TABLAS		[]
%%%%%%%%%%%%%%%%%%%%%%%%%%%%%%%%%%

%%%%%%%%%%%%%%%%%%%%%%%%%%%%%%%%%%
%(8) ÍNDICE DE ILUSTRACIONES	[]
\newpage
\listoffigures
\newrgbcolor{xdxdff}{0.49 0.49 1}
\newrgbcolor{ffqqtt}{1 0 0.2}
\psset{xunit=1.0cm,yunit=1.0cm,algebraic=true,dotstyle=*,dotsize=3pt
0,linewidth=0.8pt,arrowsize=3pt 2,arrowinset=0.25}
%%%%%%%%%%%%%%%%%%%%%%%%%%%%%%%%%%

%%%%%%%%%%%%%%%%%%%%%%%%%%%%%%%%%%
%(9) RESUMEN			[]
%%%%%%%%%%%%%%%%%%%%%%%%%%%%%%%%%%
% \newpage
% \chapter*{Resumen}
% \addcontentsline{toc}{chapter}{Resumen} % si queremos que aparezca en el índice
% Una bonita historia.


%TEXTO
%(1) INTRODUCCIÓN	[]
%(2) CUERPO DE LA TESIS	[]
%(3) CONCLUSIONES	[]
% 

\newpage
\pagenumbering{arabic}
\include{Introduccion}		%Introducción
\include{Capitulo_01}		%Introducción histórica del Infinito
\include{Capitulo_02}		%Marco Teórico
\include{Capitulo_03}		%Estudio matemático del infinito
\include{Capitulo_04}		%Plataforma virtual
\include{Conclusiones}		%Resultados, conclusiones y recomendaciones
\include{Anexo}			%Recursos educativos y otros


%PÁGINAS FINALES

%(1) GLOSARIO (Índice Analítico)	[]
\newpage
\printindex
\addcontentsline{toc}{chapter}{Índice Analítico}	%Para que aparezca en el índice

%(2) BIBLIOGRAFÍA			[]
\newpage
\bibliography{bibliografia}
\bibliographystyle{flexbib}
% \nocite{KasnerNewman, Martinez, Niven, Kaplun, Socioepistemologia, Hitt2007} Tener referencia en la bibliografía que no ha sido citada.

%(3) ANEXOS(OPTATIVA)			[]
%(4) APÉNDICE(OPTATIVA)			[]
%(5) MATERIAL COMPLEMENTARIO(OPTATIVO)	[]
\end{document}